\section{What the Plan Cache Keys On}
\label{sec:ch17-what-the-plan-cache-keys-on}
\index{plan cache!what it keys on|(}%

The cache itself is one line of declaration, a \code{ristretto} cache from a
\code{uint64} to a prepared plan. What matters is the key, and the key is the
number section~\ref{sec:ch17-what-normalization-makes} just watched being made:

\begin{minted}[fontsize=\footnotesize,breaklines]{go}
	operationID := opContext.internalHash
	// try to get a prepared plan for this operation ID from the cache
	cachedPlan, ok := p.planCache.Get(operationID)
\end{minted}

Planning on a miss is single-flighted on the same identifier, so a burst of
identical first-time requests plans once and the rest wait. That is the same
arrangement chapter~\ref{ch:inside-hotchocolate} found around compilation in
HotChocolate, arrived at independently by a different team in a different
language, which is either convergent design or a sign that everybody eventually
meets the same thundering herd.

\subsection*{Measuring it without a metrics endpoint}

\index{plan cache!why no earlier chapter could see it}%
The reason four chapters left this open is that a stock router will not tell
you. It registers cache metrics for seven caches by name and publishes none of
them, because the Prometheus setting that turns them on defaults to false.
Chapter~\ref{ch:enter-the-router} saw the metrics endpoint in the startup log
and there was nothing on it worth reading.

There is a better route than the counters anyway, and it is exact rather than
sampled. The router will evaluate expressions into access-log fields, and one of
the expressions is the plan cache key itself:

\begin{minted}[fontsize=\footnotesize,breaklines]{yaml}
access_logs:
  enabled: true
  router:
    fields:
      - key: planHash
        value_from:
          expression: request.operation.queryPlanHash
      - key: planHit
        value_from:
          expression: request.operation.planCacheHit
\end{minted}

That is a router reporting its own cache key, per request, with no counter in
between. Everything below is read off that log.

\subsection*{Eleven documents, one plan}

\index{plan cache!what it ignores}%
Eleven documents, each sent twice, land on one key:

\begin{minted}[fontsize=\footnotesize,breaklines]{text}
baseline     plan=17972416575696450974
named-a      plan=17972416575696450974
named-b      plan=17972416575696450974
var-n        plan=17972416575696450974
var-count    plan=17972416575696450974
literal-7    plan=17972416575696450974
declared     plan=17972416575696450974
declared-9   plan=17972416575696450974
spaced       plan=17972416575696450974
commented    plan=17972416575696450974
fragment     plan=17972416575696450974
\end{minted}

Reading that list against
section~\ref{sec:ch17-what-normalization-makes}: the operation name is blanked,
so \code{query A} and \code{query B} collapse. Variables are renamed
positionally, so \code{\$n} and \code{\$count} collapse. Literals are lifted
into the side channel, so \code{first: 3} and \code{first: 7} collapse, and so
does a properly declared variable carrying either value. Whitespace and comments
never survive a parse. Fragments are inlined at the spread's position, so a
fragment and the same selection written out collapse too.

\index{plan cache!the control}%
One shared plan answering two different questions is either the point or a bug,
so I checked rather than assuming:

\begin{minted}[fontsize=\footnotesize,breaklines]{text}
control: first: 3 -> 3 nodes, first: 7 -> 7 nodes
\end{minted}

The plan is the shape of the work and the values ride in beside it. So a client
that inlines its arguments costs the router nothing, which is a friendlier
answer than I expected. Chapter~\ref{ch:the-life-of-a-request} argued for
putting identifiers in the Variables pane instead, and that argument is
untouched: it was about a subgraph's operation cache, which is a different cache
in a different process, and it still holds there.

\subsection*{Three things that buy you a second plan}

\index{plan cache!what splits it}%
Now the half I did not expect. Three pairs, each differing in exactly one thing:

\begin{minted}[fontsize=\footnotesize,breaklines]{text}
alias        varies the alias on a root field
             12452750417983769066 vs 2728077130562092255
order        varies the order of two fields in a selection set
             15843417663884115703 vs 1777378001281248594
nullability  varies the declared type of a variable
             17972416575696450974 vs 16573561200335999562
\end{minted}

Those twenty-digit numbers are properties of this composed graph rather than
facts worth remembering. Recompose and they all change; what does not change is
which of them match.

Put that beside the previous list and the shape of it is odd. Asking for a
completely different slice of the catalogue is free. Renaming a field in your
own response is not. Writing \code{\{ title sku \}} where somebody else wrote
\code{\{ sku title \}} is not, because nothing normalises the order and the two
print differently. And adding a \code{!} to a variable you declared is not,
because the declared type survives into the printed document where the lifted
literal's value did not.

\index{plan cache!field order and generated clients}%
The field-order one is the one I would worry about in a real deployment, and it
is invisible from every artefact anybody looks at. A client that assembles
selection sets from a set, or a map, or anything else without a defined
iteration order will emit the same query in several orders across processes and
restarts, and each order is a separate plan occupying a separate entry in a
cache that holds 1024 of them by default. Nothing warns you. The queries are
identical, the answers are identical, and one query quietly occupies as many
entries as its fields have orderings.

\subsection*{Chapter~\ref{ch:enter-the-router} was right, and this is the rest of it}

\index{plan cache!the nullability split}%
Chapter~\ref{ch:enter-the-router} found that a properly declared variable
produces a different plan \emph{document} from the inlined form, and was careful
about what that licensed: it said outright that this showed two documents
differing, not a cache holding two entries, and handed the cache question here.
Both halves are true and the reason they look contradictory is the type. Here
are the normalised documents for five forms:

\begin{minted}[fontsize=\scriptsize,breaklines]{text}
inline-3       query($a: Int){browseProducts(first: $a){nodes {title}}}
declared-null  query($a: Int){browseProducts(first: $a){nodes {title}}}
declared-nonnl query($a: Int!){browseProducts(first: $a){nodes {title}}}
named+nonnull  query Storefront($a: Int!){browseProducts(first: $a){nodes {title}}}
named+null     query Storefront($a: Int){browseProducts(first: $a){nodes {title}}}
\end{minted}

The extraction pass synthesises a nullable variable, because a lifted literal
cannot be required. Chapter~\ref{ch:enter-the-router}'s declared form was
\code{\$first: Int!}, so its document differed by that one character and it was
looking at a genuine difference. Declare the variable nullable and the document
is identical to the inlined one, and so is the key.

\index{plan cache!two documents, one entry}%
The last two lines are the awkward part. Those two documents differ, visibly, in
the operation name, and they share a cache key, because the name was blanked
before the hash and not before the print. Whichever arrives first is planned,
stored, and served to the other. The plan a request executes is correct either
way, since execution runs the fetch tree and the fetch tree is the same; what
differs is the \code{normalizedQuery} string carried along with it, which now
names an operation the second caller did not send.

I want to be careful about how much that is worth. It is cosmetic as far as I
can tell, and I could not construct a case where it changes an answer. It is
also permanently unobservable from outside, for a reason
section~\ref{sec:ch17-watching-it-changes-it} gets to: the only way to make the
router show you a plan document is a header that stops it using the cache. So
the mismatch exists exactly where nobody can look at it.

\medskip
That is which plan you get. The next question is what a plan is, and how the
router decided that this subgraph rather than that one should answer a field
both of them offer.

\index{plan cache!what it keys on|)}%
